% ============================================================================
% Seccion 4: Implementacion de la Aplicacion
% TFP - Rubrica: 3 puntos
% ============================================================================
\section{Implementaci\'on de la aplicaci\'on}
\label{sec:implementation}

El sistema desacopla la interpretaci\'on del lenguaje (agente LLM) de la
decisi\'on de ranking (motor de scoring determin\'istico), lo que garantiza
transparencia y reproducibilidad: el LLM interpreta y explica, pero no decide el
orden de las recomendaciones.

Para complementar esta descripci\'on, se incluye la grabaci\'on del video del
proyecto en el siguiente enlace:
\begin{itemize}[noitemsep]
  \item \url{https://youtu.be/c_wMR6HVe1M}
\end{itemize}

\subsection{Arquitectura general de la soluci\'on}

El flujo principal de la soluci\'on (Figura~\ref{fig:architecture}) es el siguiente:

\begin{enumerate}[noitemsep]
  \item El usuario expresa sus preferencias en \emph{lenguaje natural} desde el
        frontend Angular.
  \item El agente conversacional descubre el perfil vocacional (RIASEC) y las
        prioridades del estudiante. El perfil y las preferencias se
        \emph{persisten} entre turnos y entre sesiones con \texttt{langmem}
        (ver detalle en la siguiente subsecci\'on) y pueden actualizarse.
  \item El agente traduce las prioridades en \emph{pesos} din\'amicos para cada
        criterio.
  \item El motor de scoring (Python) asigna un puntaje a cada
        carrera--universidad sobre \texttt{features.csv}, aplicando adem\'as
        filtros duros (regi\'on, tipo de instituci\'on, presupuesto).
  \item Se construye un ranking Top-N con los mayores puntajes.
  \item El agente redacta una explicaci\'on personalizada citando datos
        verificables.
\end{enumerate}

\begin{figure}[H]
    \centering
    \includegraphics[width=0.95\linewidth]{fig-arquitectura-logica_v2.png}
    \caption{Arquitectura l\'ogica de componentes: el frontend se conecta
    \textbf{directamente} al backend (autenticaci\'on/CRUD) y al agente
    (conversaci\'on), sin proxy entre ambos.}
    \label{fig:architecture}
\end{figure}

\emph{Nota:} el frontend Angular llama a \textbf{ambos servicios de forma
directa} --- al backend serverless para autenticaci\'on y operaciones CRUD, y
al agente (\texttt{POST /ag-ui}, \emph{streaming} SSE) para la conversaci\'on
--- reutilizando el mismo JWT en ambas llamadas; \textbf{no hay} un
\emph{proxy} del backend hacia el agente, decisi\'on documentada
expl\'icitamente en el ADR-012 del backend. El agente corre detr\'as de su
propio \emph{Application Load Balancer} p\'ublico (no detr\'as de API
Gateway), precisamente porque el \emph{timeout} fijo de 30s de API Gateway
HTTP API cortar\'ia el \emph{streaming} SSE de \texttt{/ag-ui}. La
Figura~\ref{fig:architecture} representa el dise\~no objetivo de la
arquitectura a nivel de componentes l\'ogicos (no de despliegue): el agente ya
conecta tanto la afinidad RIASEC (cat\'alogo piloto) como los datos reales de
\texttt{features.csv} v\'ia \texttt{search\_programs}; lo que falta del nodo
marcado en rojo es el \textbf{motor de scoring} que combine ambos en un
puntaje ponderado \'unico (detalle m\'as abajo). La base de datos
relacional del proyecto es \textbf{Amazon RDS para PostgreSQL} (no Aurora); se
usa hoy para la persistencia de memoria del agente (ver subsecci\'on
siguiente) y para el dominio \texttt{identity} del backend.

\begin{landscape}
\begin{figure}[H]
    \centering
    \includegraphics[width=0.92\linewidth]{fig-secuencia-conversacion_v2.png}
    \caption{Secuencia de una conversaci\'on t\'ipica, desde el mensaje del
    estudiante hasta la respuesta con recomendaciones explicadas.}
    \label{fig:conversation-sequence}
\end{figure}
\end{landscape}

\subsection{Agente conversacional (Deep Agent / deepagents)}

El agente se implementa en \textbf{Python} como un \emph{Deep Agent} construido
con el \emph{harness} \texttt{deepagents} ---del ecosistema
\emph{LangChain/LangGraph}--- que \textbf{coordina tres subagentes
especializados}: \texttt{assessment} (perfil vocacional RIASEC),
\texttt{matching} (afinidad, ranking de carreras y su aterrizaje en
universidades e institutos reales del Per\'u) y \texttt{planning} (plan de
acci\'on para la carrera elegida; ver Secci\'on~\ref{sec:prompts} para el
detalle de cada \emph{system prompt}).

\paragraph{Despliegue: contenedor en ECS Fargate, no AgentCore Runtime.} El
agente corri\'o inicialmente sobre AWS Bedrock AgentCore; la infraestructura
actual (m\'odulo \texttt{agent-service} de
\texttt{spark-match-02-infrastructure}) lo empaqueta en un \textbf{contenedor
Docker} y lo despliega en \textbf{AWS ECS Fargate} (ARM64) detr\'as de un
Application Load Balancer (raz\'on del cambio y detalle de despliegue en la
Secci\'on~\ref{sec:orchestration}).

\paragraph{Memoria del agente.} El perfil y las preferencias del estudiante se
persisten con \texttt{langmem} (\texttt{create\_manage\_memory\_tool} /
\texttt{create\_search\_memory\_tool}) bajo un espacio de nombres asociado al
usuario autenticado (JWT), respaldado por \textbf{RDS PostgreSQL} en producci\'on
(en memoria durante desarrollo/tests). En paralelo, un extractor de perfil en
segundo plano (\texttt{StudentProfile}) resume la conversaci\'on completa de
forma as\'incrona. Antes de persistir cualquier contenido, se redactan correos,
DNI y tel\'efonos peruanos (\emph{PII redaction}), de modo que la memoria de
largo plazo nunca almacena datos personales sensibles en texto plano.

El LLM (Claude en Bedrock) se resuelve a partir de la configuraci\'on de la
aplicaci\'on, que define dos modelos --- uno robusto para razonamiento y uno
r\'apido/econ\'omico para turnos simples, seleccionados por el enrutador de
intenci\'on descrito en la Secci\'on~\ref{sec:model-data} ---:

\begin{lstlisting}[style=pythonstyle, caption={Configuracion real de modelos (src/config/settings.py)}, label={lst:agent}]
model_id: str = "us.anthropic.claude-sonnet-4-5-20250929-v1:0"       # razonamiento
fast_model_id: str = "us.anthropic.claude-haiku-4-5-20251001-v1:0"   # turnos simples
\end{lstlisting}

\begin{figure}[H]
    \centering
    \includegraphics[width=\linewidth]{fig-estados-agente_v2.png}
    \caption{M\'aquina de estados del Deep Agent: enrutamiento por intenci\'on
    y coordinaci\'on de los tres subagentes especializados.}
    \label{fig:agent-states}
\end{figure}

\subsection{Motor de scoring (5 criterios)}

Cuando el agente entrega los pesos, el motor asigna un puntaje a cada
combinaci\'on carrera--universidad. La f\'ormula, con cinco criterios, es:

\[
\text{Score} = w_1\,\text{Afinidad} + w_2\,\text{Ingreso} + w_3\,\text{Costo}
             + w_4\,\text{Admisi\'on} + w_5\,\text{Duraci\'on}
\]

donde la \textbf{Afinidad} se calcula con el modelo \emph{RIASEC} de Holland
(compatibilidad entre el c\'odigo de 3 letras del estudiante y el de la carrera);
el \textbf{Ingreso} usa el sueldo promedio de egresados; el \textbf{Costo}
favorece las opciones m\'as accesibles; la \textbf{Admisi\'on} se basa en la tasa
de ingresantes/postulantes; y la \textbf{Duraci\'on} favorece carreras m\'as
cortas. Todas las variables vienen normalizadas en $[0,1]$ desde el pipeline. Los
pesos $w_1,\dots,w_5$ \textbf{no son fijos}: los infiere din\'amicamente el LLM en
cada sesi\'on, suman 1 y pueden ser 0 si el estudiante indica que un criterio no
le importa.

\begin{table}[H]
\centering
\begin{tabular}{@{}lcccccc@{}}
\toprule
\textbf{Carrera -- Universidad} & \textbf{Afin.} & \textbf{Ing.} & \textbf{Costo} & \textbf{Adm.} & \textbf{Dur.} & \textbf{Score} \\
\midrule
Estad\'istica -- UNMSM     & 95 & 85 & 90 & 60 & 80 & \textbf{86.0} \\
Econom\'ia -- UNMSM        & 80 & 82 & 88 & 55 & 80 & 79.2 \\
Ing.\ Industrial -- UNI    & 70 & 80 & 70 & 50 & 60 & 69.0 \\
\bottomrule
\end{tabular}
\caption{Ejemplo ilustrativo con pesos $(0.35, 0.25, 0.15, 0.10, 0.15)$. Los valores se muestran en escala 0--100 para facilitar la lectura; el pipeline trabaja en $[0,1]$.}
\label{tab:scoring-example}
\end{table}

Como ilustra la Tabla~\ref{tab:scoring-example}, el sistema ordena los resultados de mayor a menor y genera un ranking Top-N.
Luego el agente redacta la justificaci\'on, por ejemplo: \textit{``Te
recomendamos Estad\'istica en la UNMSM porque presenta una fuerte alineaci\'on
con tus intereses cuantitativos, tiene uno de los mejores ingresos promedio
dentro de las carreras anal\'iticas y mantiene costos relativamente bajos''.}

\paragraph{Estado de implementaci\'on.} La herramienta
\texttt{calculate\_affinity} del agente ya calcula el criterio de
\textbf{afinidad} (similitud RIASEC posicional) y ordena un Top-N por ese
criterio --- es, a la fecha, el \'unico de los cinco que el motor de
\emph{scoring} pondera dentro de un puntaje. Los otros cuatro criterios ---
\textbf{ingreso}, \textbf{costo}, \textbf{admisi\'on} y \textbf{duraci\'on} ---
ya est\'an calculados y normalizados por el \emph{pipeline} de datos
(Secci\'on~\ref{sec:model-data}), y desde la herramienta
\texttt{search\_programs} el agente puede consultarlos directamente sobre el
dataset completo de 6\,208 combinaciones carrera--instituci\'on, con cifras
reales de costo, ingreso, admisi\'on y duraci\'on en vez de datos gen\'ericos.
Lo que todav\'ia falta es la \textbf{funci\'on de combinaci\'on}: hoy el agente
presenta esos cuatro criterios junto a la afinidad para que el LLM razone
sobre ellos en su respuesta, pero no existe a\'un el c\'odigo que los combine
en un \'unico puntaje ponderado con los pesos $w_1,\dots,w_5$ de la
f\'ormula anterior; el ejemplo num\'erico de la Tabla~\ref{tab:scoring-example}
sigue siendo, por tanto, ilustrativo del dise\~no objetivo y no un resultado
medido. Ver Secci\'on~\ref{sec:results} para el detalle del alcance
implementado y el trabajo pendiente.

\subsection{Backend serverless}

El backend de plataforma (operaciones CRUD, gesti\'on de sesiones) sigue un
enfoque \emph{serverless} orientado a dominios (DDD): \textbf{TypeScript} sobre
AWS Lambda + EventBridge, empaquetado con AWS SAM. A la fecha implementa
\'unicamente el contexto \texttt{identity} (registro, login, roles,
desactivaci\'on de usuarios y auditor\'ia); el contexto de \emph{matching}/
recomendaci\'on todav\'ia no existe en esta capa (ver
Secci\'on~\ref{sec:results}). El backend cumple el rol de exposici\'on y
operaciones CRUD, sin l\'ogica de IA.

\subsection{Frontend}

El frontend es una SPA en \textbf{Angular 21} con Angular Material, organizada
por \emph{features} (\texttt{auth}, \texttt{chat}, \texttt{assessment},
\texttt{careers}, \texttt{filters}, \texttt{results}, \texttt{reports},
\texttt{profile} y \texttt{landing}). Cada feature tiene su propio
componente, servicio y pruebas unitarias; ya no es un simple \emph{scaffold}
vac\'io, aunque su integraci\'on end-to-end con el agente desplegado todav\'ia
est\'a en desarrollo.

Para acceder a la app es necesario crear una cuenta llenando el formulario de
registro con nombre completo, correo, contrase\~na, edad, regi\'on y \emph{área
de interés}. Luego se ingresa con el correo y la contrase\~na creados. En la
demo de prueba se puede usar el usuario \texttt{ftapara@unsa.edu.pe} con
contrase\~na \texttt{qwerty123}. El registro se realiza en
\url{https://d3sx4v4l1enpvk.cloudfront.net/auth/register}.

\begin{figure}[H]
    \centering
    \includegraphics[width=0.95\linewidth]{crear cuenta.png}
    \caption{Formulario de registro del frontend de Spark Match.}
    \label{fig:fw-register}
\end{figure}

\begin{figure}[H]
    \centering
    
\includegraphics[width=0.95\linewidth]{iniciar sesion.png}
    \caption{Pantalla de inicio de sesi\'on del frontend.}
    \label{fig:fw-login}
\end{figure}

\begin{figure}[H]
    \centering
    
\includegraphics[width=0.95\linewidth]{inicio.png}
    \caption{Vista de inicio del usuario en la aplicaci\'on Spark Match.}
    \label{fig:fw-home}
\end{figure}

\begin{figure}[H]
    \centering
    \includegraphics[width=0.95\linewidth]{perfil.png}
    \caption{Pantalla de perfil del estudiante, donde se muestran sus datos y preferencias.}
    \label{fig:fw-profile}
\end{figure}

\begin{figure}[H]
    \centering
    \includegraphics[width=0.95\linewidth]{filtros.png}
    \caption{Vista de filtros del frontend, que permite ajustar la orientaci\'on por regi\'on, tipo de instituci\'on y presupuesto.}
    \label{fig:fw-filters}
\end{figure}

La grabaci\'on del demo est\'a disponible en
\url{https://youtu.be/c_wMR6HVe1M} y muestra el flujo completo de registro,
inicio de sesi\'on, entrevista vocacional y recomendaciones explicadas.

El dise\~no privilegia la separaci\'on de responsabilidades (UI / API /
agente / scoring / datos), la configuraci\'on por variables de entorno y una
API sin estado (las sesiones se persisten en base de datos, no en memoria).

\paragraph{Lenguaje del n\'ucleo.} El \textbf{n\'ucleo de IA/ML del proyecto est\'a
implementado en Python}: el agente conversacional (\texttt{deepagents}), el motor
de scoring (Pandas/NumPy), el \emph{pipeline} de datos y el framework de
evaluaci\'on.
